% !PRINTED 371
\weberpage{371}
% Source-faithful three-stroke Cayley configuration sign.
\providecommand{\weberthreeup}{\mathord{\tikz[baseline=-0.2ex,x=0.23em,y=0.23em,line width=0.45pt]{\draw (0,0)--(-1,1) (0,0)--(0,1.08) (0,0)--(1,1);}}}
Es soll jetzt zunächst untersucht werden, wie sich in dieser
Bezeichnungsweise die azygetischen und die syzygetischen Tripel
unterscheiden.

Dabei ist zu beachten, dass nach der bis jetzt gegebenen
Erklärung die Ziffer 8 eine besondere Stelle einnimmt, während
die Ziffern 1 bis 7 vollständig gleichartig auftreten und beliebig
permutirt werden können.

Wir leiten aus den beiden Complexen (1), (2) noch die Ergänzung
 zu einem azygetischen Complextripel her, nämlich:
\weq{3}{
\begin{aligned}
&x_2 y_2,\ x_3 y_3,\ x_4 y_4,\ x_5 y_5,\ x_6 y_6,\ x_7 y_7,\\
&x_1 y_2,\ x_3 z_3,\ x_4 z_4,\ x_5 z_5,\ x_6 z_6,\ x_7 z_7,\\
&x_1 x_2,\ y_3 z_3,\ y_4 z_4,\ y_5 z_5,\ y_6 z_6,\ y_7 z_7,
\end{aligned}
}
und gehen nun die einzelnen Zeichen für die Doppeltangententripel
 durch.

Wir beginnen mit dem Zeichen $\weberthreeup$, für welches wir mit
Rücksicht auf die Ausnahmestellung der Ziffer 8 drei Typen zu
betrachten haben:
\weq{4}{
\begin{aligned}
&[1\,8]\,[2\,8]\,[3\,8] = x_1 x_2 x_3\\
&[1\,5]\,[1\,6]\,[1\,7] = y_5 y_6 y_7\\
&[1\,6]\,[1\,7]\,[1\,8] = y_6 y_7 x_1,
\end{aligned}
}
und der Anblick der drei Complexe (3) zeigt (nach dem Satze
\S.\,96, 3.), dass alle diese Tripel azygetisch sind.

Zweitens betrachten wir das Zeichen $\vee\mid$, für welches vier
Typen zu berücksichtigen sind:
\weq{5}{
\begin{aligned}
&[8\,3]\,[8\,4]\,[1\,2] = x_3 x_4 y_2\\
&[8\,4]\,[1\,4]\,[2\,3] = x_4 y_4 z_3\\
&[1\,2]\,[1\,3]\,[4\,8] = y_2 y_3 x_4\\
&[1\,4]\,[1\,5]\,[2\,3] = y_4 y_5 z_3,
\end{aligned}
}
und auch diese Tripel sind azygetisch.

Für das Zeichen $\triangle$ ist zu betrachten:
\weq{6}{
\begin{aligned}
&[1\,8]\,[2\,8]\,[1\,2] = x_1 x_2 y_2\\
&[1\,2]\,[1\,3]\,[2\,3] = y_2 y_3 z_3,
\end{aligned}
}
die sich gleichfalls als azygetisch erweisen, weil $y_2$ in dem Complexe
 (3), der die Paare $x_1 x_2$, $y_3 z_3$ enthält, nicht vorkommt.

Für das Zeichen $\mid\mid\mid$ giebt es zwei Typen:
\weq{7}{
\begin{aligned}
&[1\,3]\,[2\,4]\,[5\,8] = y_3 z_4 x_5\\
&[1\,3]\,[2\,4]\,[5\,6] = y_3 z_4 [5\,6].
\end{aligned}
}
